\documentclass[11pt,a4paper]{article}
\usepackage[T1]{fontenc}
\usepackage[utf8]{inputenc}
\usepackage{lmodern}
\usepackage[margin=27mm,headheight=15pt]{geometry}
\usepackage{amsmath,amssymb,amsthm,mathtools}
\usepackage{booktabs,longtable,array,enumitem,microtype}
\usepackage{xcolor,fancyhdr}
\usepackage[unicode,colorlinks=true,linkcolor=blue!45!black,citecolor=blue!45!black,urlcolor=blue!45!black]{hyperref}
\usepackage{bookmark}
\definecolor{ink}{RGB}{25,42,64}
\definecolor{accent}{RGB}{35,91,120}
\newcommand{\R}{\mathbb R}
\newcommand{\N}{\mathbb N}
\newcommand{\E}{\mathbb E}
\newcommand{\Pbb}{\mathbb P}
\newcommand{\norm}[1]{\left\lVert#1\right\rVert}
\newcommand{\abs}[1]{\left\lvert#1\right\rvert}
\newcommand{\dd}{\,\mathrm d}
\DeclareMathOperator{\dist}{dist}
\DeclareMathOperator{\supp}{supp}
\DeclareMathOperator{\spanop}{span}
\DeclareMathOperator{\diam}{diam}
\DeclareMathOperator{\Lip}{Lip}
\newtheorem{theorem}{Theorem}[section]
\newtheorem{lemma}[theorem]{Lemma}
\newtheorem{proposition}[theorem]{Proposition}
\newtheorem{corollary}[theorem]{Corollary}
\theoremstyle{definition}
\newtheorem{definition}[theorem]{Definition}
\theoremstyle{remark}
\newtheorem{remark}[theorem]{Remark}
\numberwithin{equation}{section}
\setlist{topsep=4pt,itemsep=2pt,parsep=0pt}
\setlength{\emergencystretch}{2em}
\pagestyle{fancy}
\fancyhf{}
\fancyhead[L]{\small\color{ink}AIM 641}
\fancyhead[R]{\small 4 November 2026}
\fancyfoot[C]{\thepage}
\renewcommand{\headrulewidth}{0.3pt}


\urlstyle{same}
\title{Quartic simplex interpolation: a tetrahedral Fej\'er theorem and a dimension-free fourth-power bound}
\author{Matthew J. Colbrook\thanks{The solutions were found by AI. The submitted arguments and the repository audit are AI work; independent human review and formal verification are not claimed.}\\[4pt]
\small Department of Applied Mathematics and Theoretical Physics\\
\small University of Cambridge, Cambridge CB3 0WA, United Kingdom\\
\small\href{mailto:m.colbrook@damtp.cam.ac.uk}{m.colbrook@damtp.cam.ac.uk}}
\date{4 November 2026}
\hypersetup{pdftitle={Quartic simplex interpolation: a tetrahedral Fej\'er theorem and a dimension-free fourth-power bound},pdfauthor={Matthew J. Colbrook},pdfsubject={AIM 641: AI-found solution or partial result}}
\begin{document}
\maketitle
\begin{abstract}
For Bos's prescribed quartic interpolation nodes, we prove the Fej\'er inequality on a tetrahedron and the fourth-power inequality in every dimension, including the equality set. Finite exact-arithmetic witnesses and reconstructing verifiers are supplied. These resolve assertions 2 and 4 of AIM 641; the all-dimensional quartic Fekete assertion remains unresolved.
\end{abstract}
\noindent\textbf{Catalogue target.} \href{https://github.com/MColbrook/AIM/blob/c4036871b0bf15cbc83e79c780b7191bcbb2cfed/problems/641-bos-simplex-interpolation-nodes.md}{AIM 641 at repository revision \texttt{c4036871b0bf}}. The package preserves its full statement in \texttt{statement.md}; the review and target comparison are linked from the package README.
\medskip

\section{Quartic simplex interpolation: Problems 641(2) and 641(4)}
\label{sec:simplex}

We prove two assertions from Problem 641: the prescribed quartic nodes are
Fej\'er points on a tetrahedron, and their cardinal polynomials satisfy the
fourth-power inequality in every dimension.  Both proofs reduce the required
inequalities to explicit finite calculations in an algebraic number field.
The accompanying standard-library Python verifiers reconstruct these
calculations exactly.  The known planar quartic Fej\'er theorem
\cite[Proposition 4.1]{bos2023}, attributed there to \cite{bos1983}, is the
base case; we supply an independent elementary proof below, making the
argument self-contained.  The all-dimensional quartic Fekete assertion
remains open in this report.

\subsection{Nodes and homogeneous cardinal polynomials}

Write $m=d+1$ and use barycentric coordinates on
\[
 \Delta_{m-1}=\{x\in[0,\infty)^m:p_1(x)=1\},
 \qquad p_j(x)=\sum_{i=1}^m x_i^j.
\]
Set
\[
 r=\sqrt5,\qquad q=\sqrt{3/7},\qquad
 a=\frac{4+r}{11},\qquad b=\frac{7-r}{22}.
\]
The node set $\mathcal F_{4,m-1}$ consists of all distinct coordinate
permutations of the following vectors, padded with zero coordinates:
\begin{equation}\label{eq:simplex-nodes}
 (1),\quad (\tfrac12,\tfrac12),\quad
 (\tfrac{1+q}{2},\tfrac{1-q}{2}),\quad
 (a,b,b),\quad (\tfrac14,\tfrac14,\tfrac14,\tfrac14).
\end{equation}
An orbit is omitted when it has more than $m$ nonzero coordinates.  The
number of nodes is
\[
 N_m=m+3\binom m2+3\binom m3+\binom m4=\binom{m+3}{4},
\]
the dimension of polynomials of degree at most four on $\Delta_{m-1}$.

We use the following degree-four homogeneous polynomials on $\R^m$.
Their restrictions to $p_1=1$ are the cardinal polynomials.  The vertex
polynomials are
\begin{equation}\label{eq:simplex-vertex}
\begin{split}
 V_i=2x_i\bigg[&\frac{101+17r}{11}x_i^3-(13+2r)x_i^2p_1
      +\frac{130+27r}{22}x_ip_1^2-\frac{259+51r}{264}p_1^3\\
 &+p_2\left(\frac{9+r}{8}p_1-\frac{24+17r}{22}x_i\right)
      +p_3\left(\frac{45+3r}{44}-\frac53\right)\bigg].
\end{split}
\end{equation}
For each unordered pair $\{i,j\}$, define
\begin{align}
 U_{ij}={}&\frac{61+7r}{22}(x_i^2+x_j^2)
       +\frac{3+9r}{11}x_ix_j
       -\frac{41+13r}{22}(x_i+x_j)p_1 \notag\\
       &+\frac{8+2r}{11}p_1^2-\frac{7-r}{11}p_2,
       \label{eq:simplex-U}\\
 W_{ij}={}&\frac{7+r}{2}(x_j^2-x_i^2)
       +\frac{5+r}{2}(x_i-x_j)p_1,\label{eq:simplex-W}\\
 E_{ij}^{\pm}={}&\frac{49}{6}x_ix_j(U_{ij}\pm qW_{ij}),
       \label{eq:simplex-edge}\\
 M_{ij}={}&\frac4{33}x_ix_j\big[(100-8r)(x_i^2+x_j^2)
       +(672-96r)x_ix_j \notag\\
       &\hspace{24mm}+(62r-302)(x_i+x_j)p_1
       +(82-40r)p_2+(76-14r)p_1^2\big].
       \label{eq:simplex-midpoint}
\end{align}
Interchanging $i,j$ interchanges the two labels $+$ and $-$.  For each
unordered triple there are three polynomials, one for each distinguished
coordinate:
\begin{equation}\label{eq:simplex-triple}
 T_{i;jk}=\frac{73+r}{2}x_ix_jx_k
 \left[(1+r)x_i+\frac{3-r}{2}(x_j+x_k)-p_1\right].
\end{equation}
For each unordered quadruple there is one polynomial,
\begin{equation}\label{eq:simplex-quadruple}
 C_{ijkl}=256x_ix_jx_kx_l.
\end{equation}
Let $\mathcal L_m$ denote this family of $N_m$ polynomials.

\begin{lemma}[Exact cardinal verification]\label{lem:simplex-cardinal}
The evaluation matrix of $\mathcal L_m$ at
$\mathcal F_{4,m-1}$ is a permutation matrix.  Consequently
$\mathcal F_{4,m-1}$ is unisolvent, and the restrictions of
$\mathcal L_m$ to $\Delta_{m-1}$ are its cardinal polynomials.
\end{lemma}

\begin{proof}
For $m=4$, direct substitution in
\eqref{eq:simplex-vertex}--\eqref{eq:simplex-quadruple} gives all
$35^2=1225$ entries exactly.  This calculation is reproduced by
\texttt{verify\_641\_exact.py}, using rational arithmetic in
$\mathbb Q(r,q)$.  There are $35$ ones, one in each row and column,
and all other entries are zero.

This finite check proves the assertion for arbitrary $m$.  Every node has
at most four nonzero coordinates.  Every polynomial involving a
distinguished coordinate that is zero at the node vanishes there.  The
remaining polynomial formulas depend only on their distinguished
coordinates and $p_1,p_2,p_3$, and so agree exactly with their
four-coordinate versions after zero coordinates are added or deleted.
Thus every possible evaluation is already covered by the $m=4$ check.
The node count is the dimension of the polynomial space, which completes
the assertion.
\end{proof}

\begin{remark}[A correction verified by interpolation]
The displayed midpoint formula in the public 2022 preprint of
\cite{bos2023} omits the term $(76-14\sqrt5)p_1^2$ in
\eqref{eq:simplex-midpoint}; without it, the midpoint value is
$(14\sqrt5-43)/33$ instead of $1$.  A paired-edge formula also repeats
one coordinate index.  We use the formulas above, whose values are checked
independently at every node.  No claim is made here about whether these
typographical issues persist in the final journal version.
\end{remark}

\subsection{An independent proof of the known planar base case}

The planar result is already known and is not counted as a new resolution.
The following short proof removes any need to assume its original
algebraic certificate.

\begin{lemma}[Planar quartic Fej\'er inequality]\label{lem:planar-fejer}
For every $x\in[0,\infty)^3$,
\[
 F_3(x):=(x_1+x_2+x_3)^8-\sum_{L\in\mathcal L_3}L(x)^2\ge0.
\]
\end{lemma}

\begin{proof}
Let $e_1=x_1+x_2+x_3$, $e_2=x_1x_2+x_1x_3+x_2x_3$ and
$e_3=x_1x_2x_3$.  Direct expansion of the cardinal formulas gives
\begin{equation}\label{eq:planar-elementary}
\begin{split}
 F_3={}&20e_1^6e_2-360e_1^4e_2^2+2100e_1^2e_2^3-3920e_2^4\\
 &+e_3\left[\frac{4623-57r}{11}e_1^5
       +\frac{-125333+7219r}{33}e_1^3e_2
       +\frac{29764-2772r}{3}e_1e_2^2\right]\\
 &+e_3^2\left[\frac{-136228+14134r}{11}e_1^2
       +\frac{77200-3296r}{3}e_2\right].
\end{split}
\end{equation}
For nonnegative coordinates, $e_2\le e_1^2/3$.  The coefficient of
$e_2$ in the last bracket is positive, so that bracket is at most
\[
 \frac{-376852+90950r}{99}e_1^2<0
 \qquad(e_1>0).
\]
Thus $F_3$ is concave as a function of $e_3$ when $e_1,e_2$ are fixed.

For fixed feasible $e_1,e_2$, the set of nonnegative coordinate triples
is compact.  At either extreme value of their product $e_3$, at least
one coordinate is zero or two coordinates are equal.  To see this,
suppose all coordinates are positive and distinct.  The gradients of
$e_1,e_2$ are then independent, so the Lagrange multiplier equations
at a product extremum have the form
$x_2x_3=\lambda+\mu(x_2+x_3)$ and its cyclic analogues.  Subtracting
two equations forces the remaining coordinate to equal $\mu$;
distinctness would force all three coordinates to equal $\mu$, a
contradiction.  Degeneracy of the constraint gradients already means
all three coordinates are equal.  Concavity therefore reduces the
desired inequality to a zero coordinate and to two equal coordinates.

Both restrictions have explicit factorizations:
\begin{align}
 F_3(u,v,0)&=20uv(u-v)^2(u^2-5uv+v^2)^2,
       \label{eq:planar-edge}\\
 F_3(u,v,v)&=40uv\left(u-\frac{3+r}{2}v\right)^2Q(u,v),
       \label{eq:planar-repeated}
\end{align}
where
\begin{equation}\label{eq:planar-quartic-sos}
 Q(u,v)=(u^2+\kappa uv+v^2)^2
          +v^2(\alpha u^2+\beta uv+\gamma v^2)
\end{equation}
with
\begin{align*}
 \kappa&=\frac{-4397+383r}{880},&
 \alpha&=\frac{23382299+5935673r}{1161600},\\
 \beta&=\frac{-229907+88485r}{1320},&
 \gamma&=\frac{39579-17344r}{165}.
\end{align*}
Here $\alpha,\gamma>0$ and
\[
 4\alpha\gamma-\beta^2
   =\frac{-529859170431+237067014259r}{11979000}>0.
\]
All displayed signs follow from the elementary bounds
$559/250<\sqrt5<9/4$.  Hence the quadratic form in
\eqref{eq:planar-quartic-sos} is positive definite, and both
\eqref{eq:planar-edge} and \eqref{eq:planar-repeated} are nonnegative
for $u,v\ge0$.  Concavity proves the result; the origin is immediate.

Every identity and algebraic sign in this proof is checked directly
from the frozen cardinal formulas by
\texttt{verify\_641\_planar\_base.py}, using only the same exact
standard-library arithmetic.  In particular, the factorization is
verified as a polynomial identity, rather than inferred from sampled
values or a numerical semidefinite calculation.
\end{proof}

\subsection{The tetrahedral Fej\'er inequality}

\begin{theorem}[Problem 641(2)]\label{thm:tetra-fejer}
For every $x\in\Delta_3$,
\begin{equation}\label{eq:tetra-fejer}
 \sum_{L\in\mathcal L_4}L(x)^2\le1.
\end{equation}
Hence the prescribed $35$ quartic nodes are also Fekete points on a
tetrahedron.
\end{theorem}

\begin{proof}
Define the degree-eight homogeneous deficit
\[
 F(x)=p_1(x)^8-\sum_{L\in\mathcal L_4}L(x)^2.
\]
The family is invariant under coordinate permutations, so $F$ is
symmetric.  The radical $q$ cancels between the two edge terms, and all
coefficients of $F$ belong to $\mathbb Q(\sqrt5)$.
It suffices to prove $F\ge0$ on the nonnegative orthant.

Order the coordinates and write
\[
 x_1=u+v+w+t,\qquad x_2=v+w+t,\qquad
 x_3=w+t,\qquad x_4=t,
 \qquad u,v,w,t\ge0.
\]
Put
\[
 G(u,v,w,t)=F(u+v+w+t,v+w+t,w+t,t),
 \qquad
 R=G(u,v,w,t)-G(u,v,w,0).
\]
The following is the finite algebraic certificate:
\begin{equation}\label{eq:fejer-cert}
 \mathcal Q(u,v,w,t)=(u+v+w+t)^{12}R(u,v,w,t)
 \quad\text{has only nonnegative coefficients.}
\end{equation}
More precisely, $R$ has $116$ nonzero coefficients and $\mathcal Q$
has $1536$.  Every nonzero coefficient of $\mathcal Q$ is at least
$60$, and each has the form
\begin{equation}\label{eq:fejer-field}
 \frac{A_\alpha+B_\alpha\sqrt5}{363},
 \qquad A_\alpha,B_\alpha\in\mathbb Z.
\end{equation}
The coefficient of $u^{19}t$ is exactly $60$.

The attached file
\texttt{fejer\_641\_certificate\_independent.json} records these
integer pairs.  The verifier regenerates them from the cardinal formulas;
it does not assume the correctness of the stored coefficient list.
For an explicit description of the computation, expand $R$ by ordinary
polynomial convolution in $\mathbb Q(\sqrt5)$, then use the recurrence
\begin{equation}\label{eq:fejer-recurrence}
 c_\alpha^{(j+1)}=\sum_{s=1}^4c_{\alpha-e_s}^{(j)},
 \qquad 0\le j<12,
\end{equation}
with missing or negative-index entries zero and $c^{(0)}$ the
coefficients of $R$.  Every sign is checked exactly by integer
comparisons: for the mixed-sign cases of $A+B\sqrt5$, compare
$A^2$ and $5B^2$.  A separate rational enclosure check uses
\[
 \frac{2236067977499789}{10^{15}}<\sqrt5<
 \frac{2236067977499790}{10^{15}},
\]
with both endpoint inequalities verified by squaring integers.
These finite operations prove \eqref{eq:fejer-cert}.

It follows that $R\ge0$, because $u+v+w+t>0$ away from the
origin and the assertion is trivial at the origin.  When $t=0$, all
polynomials involving coordinate four vanish, and the remaining family
is exactly the planar quartic cardinal family.
Lemma~\ref{lem:planar-fejer} gives $G(u,v,w,0)\ge0$ directly in
homogeneous form.  Hence $G\ge0$, and symmetry proves $F\ge0$.

For the Fekete conclusion, evaluate the cardinal family at any competing
$35$ nodes and call the resulting matrix $\mathsf L$.  Each row has
Euclidean norm at most one by \eqref{eq:tetra-fejer}.  Hadamard's
inequality gives $|\det\mathsf L|\le1$.  The determinant of
$\mathsf L$ is the ratio of the competing Vandermonde determinant to
the one at the prescribed nodes, for any fixed polynomial basis.
Thus the latter nodes maximize the absolute Vandermonde determinant.
\end{proof}

\subsection{A stable symmetric deficit for fourth powers}

\begin{theorem}[Problem 641(4)]\label{thm:simplex-fourth}
For every $m\ge2$ and every $x\in\Delta_{m-1}$,
\begin{equation}\label{eq:simplex-fourth}
 \sum_{L\in\mathcal L_m}|L(x)|^4\le1.
\end{equation}
In particular, every individual quartic cardinal polynomial is bounded
in absolute value by one, in every dimension.
\end{theorem}

We first construct one symmetric polynomial which works for every
coordinate count.  For $x\in[0,\infty)^m$, put
\[
 H_m(x)=p_1(x)^{16}-\sum_{L\in\mathcal L_m}L(x)^4.
\]
There is a polynomial
$H\in\mathbb Q(\sqrt5)[p_1,\ldots,p_{16}]$, independent of $m$,
with weighted degree $16$, whose specialization is $H_m$.
Here the weight of $p_j$ is $j$.
It has $163$ nonzero power-sum monomials.  This polynomial can be
reconstructed from the displayed cardinal formulas using only
\begin{align}
 \sum_i x_i^a={}&p_a,\notag\\
 \sum_{i\ne j}x_i^ax_j^b={}&p_ap_b-p_{a+b},\notag\\
 \sum_{\substack{i,j,k\\\text{distinct}}}x_i^ax_j^bx_k^c
 ={}&p_ap_bp_c-p_{a+b}p_c-p_{a+c}p_b-p_{b+c}p_a
       +2p_{a+b+c},\label{eq:stable-sums}\\
 \sum_{i<j<k<l}x_i^4x_j^4x_k^4x_l^4
 ={}&\frac{p_4^4-6p_4^2p_8+3p_8^2+8p_4p_{12}-6p_{16}}{24}.
       \notag
\end{align}
The exponents $a,b,c$ in the distinguished variables are positive.
For the paired-edge and midpoint contribution, apply the ordered-pair
identity to $(E_{ij}^{+})^4+(E_{ij}^{-})^4+M_{ij}^4$ and divide by
two.  For the triple contribution, apply the ordered-triple identity to
$T_{i;jk}^4$ and divide by two, since its two nondistinguished indices
are interchangeable.  The last line treats the quadruple contribution.
As before, the radical $q$ cancels.

Define formal derivatives in the polynomial ring by
\begin{equation}\label{eq:AB}
 A=\frac{\partial H}{\partial p_1},\qquad
 B=\sum_{k=2}^{16}k p_{k-1}\frac{\partial H}{\partial p_k}.
\end{equation}
The chain rule gives
\begin{equation}\label{eq:directional-deficit}
 D_m(x):=\sum_{i=1}^m\partial_{x_i}H_m(x)
       =mA(p(x))+B(p(x)).
\end{equation}
The next certificate is dimension independent.

\begin{lemma}[Exact derivative certificate]\label{lem:derivative-cert}
For the power sums of any finite collection of nonnegative numbers,
\begin{equation}\label{eq:derivative-cert}
 p_1^8A\ge0,\qquad p_1^8(5A+B)\ge0.
\end{equation}
Consequently $D_m(x)\ge0$ for every $m\ge5$ and
$x\in[0,\infty)^m$.
\end{lemma}

\begin{proof}
For a partition $\lambda=(\lambda_1,\ldots,\lambda_\ell)$ into
positive parts, let $m_\lambda^{(m)}$ be the monomial symmetric
polynomial in $m$ variables: each distinct monomial with nonzero
exponent multiset $\lambda$ occurs once.  It is zero when $m<\ell$.
The first exact expansion is
\begin{equation}\label{eq:A-monomials}
 p_1^8A=\sum_{\lambda\vdash23}a_\lambda m_\lambda,
 \qquad a_\lambda\ge0.
\end{equation}
All $1255$ partition types of $23$ occur.  For the second polynomial,
write
\[
 p_1^8(5A+B)=\sum_{\lambda\vdash23}b_\lambda m_\lambda.
\]
Exactly eleven coefficients $b_\lambda$ are negative.  They are
removed by the $28$ same-length Muirhead transfers in
Table~\ref{tab:muirhead-payments}.

The normalization is essential.  If $\mu_j$ are the multiplicities of
the distinct positive parts of $\lambda$, put
\begin{equation}\label{eq:mean-normalization}
 o_\lambda=\frac{\ell!}{\prod_j\mu_j!},\qquad
 M_\lambda^{(m)}=\frac{m_\lambda^{(m)}}{o_\lambda}.
\end{equation}
For each $\ell$-element subset of the coordinates, the corresponding
contribution to $M_\lambda^{(m)}$ is the average of all distinct
permutations of the exponent vector $\lambda$.  Therefore, if
$\lambda$ majorizes $\nu$ and they have the \emph{same length},
Muirhead's inequality\cite{muirhead} applied on each such coordinate subset yields
\begin{equation}\label{eq:muirhead-dimension-free}
 M_\lambda^{(m)}\ge M_\nu^{(m)}\qquad\text{for every }m.
\end{equation}
For $m$ smaller than their common length both sides are zero.
There is no dimension-dependent combinatorial factor in
\eqref{eq:muirhead-dimension-free}.  This would generally be false if
one transferred a normalized inequality between partitions of different
lengths without correcting their multiplicities.

The exact identity is
\begin{equation}\label{eq:muirhead-identity}
 p_1^8(5A+B)=
 \sum_{e=1}^{28}c_e\big(M_{\lambda_e}-M_{\nu_e}\big)
 +\sum_{\lambda\vdash23}r_\lambda M_\lambda,
 \qquad c_e>0,\quad r_\lambda\ge0.
\end{equation}
Every pair in the first sum has equal length and satisfies all the
majorization prefix inequalities.  There are nine length-three pairs,
eighteen length-four pairs, and one length-five pair.
The integers in Table~\ref{tab:muirhead-payments} specify
\[
 c_e=\frac{A_e+B_e\sqrt5}{395307}.
\]
The coefficients in \eqref{eq:A-monomials}, the coefficients
$b_\lambda$, and the residual coefficients $r_\lambda$ have the same
common denominator and are all supplied in
\texttt{fourth\_641\_muirhead\_certificate.json}.
Equivalently, the residuals are uniquely determined from the table by
\begin{equation}\label{eq:muirhead-residual}
 r_\lambda=o_\lambda b_\lambda
       -\sum_{e:\lambda_e=\lambda}c_e
       +\sum_{e:\nu_e=\lambda}c_e.
\end{equation}
There are $1227$ nonzero residual coefficients.

For reproducibility, the verifier expands $H$ using
\eqref{eq:stable-sums}, differentiates to obtain $A,B$, and converts
power-sum products to monomial coefficients by the following elementary
recurrence.  If $C_\mu$ is the coefficient of $m_\mu$, the
coefficient of $m_\nu$ in $p_k\sum C_\mu m_\mu$ is
\begin{equation}\label{eq:symmetric-recurrence}
 \sum_{j:\nu_j\ge k}
 C_{\operatorname{sort}(\nu-k e_j)}.
\end{equation}
Zero parts are deleted and equal positions are counted separately.
Indeed, consider the coefficient of one fixed monomial: its factor
from $p_k$ can be supplied by any coordinate whose exponent is at least
$k$.  Repeated multiplication with $k=1$ gives the factor $p_1^8$.
The verifier checks every coefficient of
\eqref{eq:A-monomials}, every majorization relation in the table, every
positive payment, every residual sign in
\eqref{eq:muirhead-residual}, and the complete identity
\eqref{eq:muirhead-identity}.  All signs are decided by exact rational
and integer arithmetic in $\mathbb Q(\sqrt5)$.

Nonnegative monomial coefficients prove the first inequality in
\eqref{eq:derivative-cert}.  Muirhead and
\eqref{eq:muirhead-identity} prove the second.  Finally,
\[
 p_1^8D_m=p_1^8(5A+B)+(m-5)p_1^8A\ge0
 \qquad(m\ge5).
\]
For nonzero $x\ge0$, division by $p_1^8>0$ is legitimate; at the
origin the result follows directly by homogeneity.
\end{proof}

\begin{proof}[Proof of Theorem~\ref{thm:simplex-fourth}]
Theorem~\ref{thm:tetra-fejer}, homogeneity, and restriction to faces
give the base cases $m\le4$:
\[
 \sum_{L\in\mathcal L_m}L(x)^4
 \le\left(\sum_{L\in\mathcal L_m}L(x)^2\right)^2
 \le p_1(x)^{16}.
\]
Suppose $m\ge5$ and the result holds for $m-1$.  Given $x\ge0$, set
\[
 a=\min_i x_i,\qquad y=x-a\boldsymbol1.
\]
At least one coordinate of $y$ is zero.  Lemma~\ref{lem:derivative-cert}
implies
\[
 H_m(x)-H_m(y)=\int_0^a D_m(y+t\boldsymbol1)\,dt\ge0.
\]
Delete a zero coordinate to obtain $\widehat y$.  The cardinal formulas
are stable under this deletion: terms with that distinguished
coordinate vanish and all other terms agree.  Hence
\[
 H_m(x)\ge H_m(y)=H_{m-1}(\widehat y)\ge0.
\]
Induction proves the homogeneous inequality in all dimensions.
Restricting to $p_1=1$ proves \eqref{eq:simplex-fourth}.
\end{proof}

\begin{corollary}[The equality set]\label{cor:simplex-equality}
For the quartic cardinal family in every dimension, equality in
\eqref{eq:simplex-fourth} holds exactly at the prescribed interpolation nodes.
\end{corollary}
\begin{proof}
The frozen certificate has $a_{(23)}=40$ and $r_{(23)}=160$.
Since $m_{(23)}=M_{(23)}=p_{23}$, the nonnegative expansions
\eqref{eq:A-monomials} and \eqref{eq:muirhead-identity} give
\[
 p_1^8D_m\ge40(m-1)p_{23}>0
 \qquad(m\ge5,\ x\ge0,\ x\ne0).
\]
If $x$ has at least five positive coordinates, first delete its zero
coordinates, obtaining a $k$-variable point with $k\ge5$ and strictly
positive minimum coordinate. The preceding minimum-coordinate
integration is then strict, so $H_k(x)>0$.

If the support has at most four coordinates, the square-sum Fej\'er
bound applies. Equality in
\[
 \sum L(x)^4\le\left(\sum L(x)^2\right)^2\le1
\]
forces exactly one cardinal value to be $1$ or $-1$ and all others to
vanish. Partition of unity forces that value to be $1$.
Reproduction of each barycentric linear coordinate then identifies
$x$ with the corresponding interpolation node. Conversely, the
cardinal equations give equality at every node.
\end{proof}

\clearpage
\subsection{Exact witness and verification scope}

The table below is the entire Muirhead transfer witness.  Its role is
to remove the eleven negative monomial coefficients of
$p_1^8(5A+B)$, leaving the nonnegative residuals given by
\eqref{eq:muirhead-residual}.  Partitions are written as lists of
positive parts, and every row has total degree $23$.

\begingroup
\scriptsize
\setlength{\tabcolsep}{3pt}
\renewcommand{\arraystretch}{1.15}
\begin{longtable}{@{}rllll@{}}
\caption{The 28 exact Muirhead payments, $c_e=(A_e+B_e\sqrt5)/395307$.}
\label{tab:muirhead-payments}\\
\toprule
$e$ & $\lambda_e$ & $\nu_e$ & $A_e$ & $B_e$ \\
\midrule
\endfirsthead
\multicolumn{5}{l}{\textit{Table \thetable\ continued}}\\
\toprule
$e$ & $\lambda_e$ & $\nu_e$ & $A_e$ & $B_e$ \\
\midrule
\endhead
\midrule
\multicolumn{5}{r}{\textit{Continued on next page}}\\
\endfoot
\bottomrule
\endlastfoot
1 & $(13,7,3)$ & $(13,6,4)$ & $1498725325337592$ & $4787655964691688$ \\
2 & $(13,7,3)$ & $(13,5,5)$ & $8536773533222352$ & $-3672905515511760$ \\
3 & $(13,8,2)$ & $(13,5,5)$ & $10632178332256140$ & $86958143772588$ \\
4 & $(14,7,2)$ & $(14,5,4)$ & $5877037193901192$ & $-79287322949208$ \\
5 & $(14,6,3)$ & $(14,5,4)$ & $1634552277469224$ & $849845887111560$ \\
6 & $(14,6,3)$ & $(13,5,5)$ & $2236877237687496$ & $-901354494250296$ \\
7 & $(13,9,1)$ & $(13,5,5)$ & $1971824645149920$ & $256388529713184$ \\
8 & $(15,6,2)$ & $(15,4,4)$ & $445045370515296$ & $-99159682090224$ \\
9 & $(15,6,2)$ & $(13,5,5)$ & $-22491224378121804$ & $10692363542691444$ \\
10 & $(8,7,5,3)$ & $(8,6,5,4)$ & $42053528057751292992$ & $323944859138228352$ \\
11 & $(8,7,5,3)$ & $(8,5,5,5)$ & $29821869380033387424$ & $-349721950787748864$ \\
12 & $(8,7,5,3)$ & $(7,7,5,4)$ & $15227149544040964512$ & $-1045096302260394240$ \\
13 & $(9,6,5,3)$ & $(7,7,5,4)$ & $34890784624009578288$ & $1842836406593505648$ \\
14 & $(9,6,5,3)$ & $(7,6,6,4)$ & $48888060142970622096$ & $-2773939678701409200$ \\
15 & $(9,7,4,3)$ & $(7,6,6,4)$ & $66937003491622708368$ & $-1294371475424784336$ \\
16 & $(8,7,6,2)$ & $(7,6,6,4)$ & $-6029137221319679280$ & $3087929711581169136$ \\
17 & $(8,7,6,2)$ & $(7,6,5,5)$ & $64746162104032292976$ & $-1885735473280680240$ \\
18 & $(10,6,4,3)$ & $(7,6,5,5)$ & $49084587467364645744$ & $-195796276763905296$ \\
19 & $(9,7,5,2)$ & $(7,6,5,5)$ & $48480702963344097792$ & $-661283746644118176$ \\
20 & $(8,6,6,3)$ & $(7,6,5,5)$ & $41670056935663240176$ & $438034863058464816$ \\
21 & $(7,7,6,3)$ & $(7,6,5,5)$ & $34233944530210993200$ & $1531159251798766608$ \\
22 & $(10,6,5,2)$ & $(7,6,5,5)$ & $4453774166762859552$ & $27706801221905184$ \\
23 & $(10,6,5,2)$ & $(6,6,6,5)$ & $33665566828687711536$ & $-1354359951228246096$ \\
24 & $(8,8,4,3)$ & $(6,6,6,5)$ & $36283080992417756064$ & $-838102952190264960$ \\
25 & $(9,6,4,4)$ & $(6,6,6,5)$ & $31927055780967653856$ & $-271911363432951264$ \\
26 & $(10,5,5,3)$ & $(6,6,6,5)$ & $28779142391558571360$ & $78345804440140896$ \\
27 & $(9,8,4,2)$ & $(6,6,6,5)$ & $-163447391863139040$ & $1812146122265513472$ \\
28 & $(8,6,5,3,1)$ & $(6,6,5,5,1)$ & $44030842344222195840$ & $-7861535751971828640$ \\
\end{longtable}
\endgroup

The verification requires Python 3 and its standard library only:
\begin{verbatim}
python3 verify_641_exact.py
python3 verify_641_planar_base.py
python3 verify_641_fourth_exact.py
\end{verbatim}
The first program checks all cardinal evaluations, symmetry and
homogeneity of the degree-eight deficit, and all $1536$ coefficient
signs.  The planar program verifies the elementary symmetric identity,
the two factorizations, concavity, and positivity of the residual
quadratic form.  The fourth-power program reconstructs the stable
polynomial $H$, checks its
$163$ terms and the exact derivative identity, and verifies
\eqref{eq:A-monomials} and \eqref{eq:muirhead-identity}.  It also
compares the power-sum expression with direct cardinal evaluations at
exact rational points with $2,3,4,5,6,9,17$ coordinates.  These pointwise
checks supplement the formal algebraic derivation; they are not used as
a substitute for it.

The two results establish 641(2) and 641(4).  The tetrahedral Fekete
property follows as a corollary, but a fourth-power bound does not imply
the Hadamard bound needed for the all-dimensional Fekete assertion
641(3).  That assertion remains unresolved here.  The cubic Fej\'er
assertion 641(1) was already marked solved in the catalogue and is not
counted as new work.  The proofs above are computer-assisted exact
proofs with explicit finite witnesses; they have not been checked by a
formal proof assistant or external human peer review.

\begin{thebibliography}{99}
\bibitem{bos2023}L. Bos, \emph{On Fekete points for a real simplex}, Indag. Math. \textbf{34} (2023), 274--293. Primary preprint: \url{https://arxiv.org/abs/2205.06498}; \url{https://doi.org/10.1016/j.indag.2022.11.003}.
\bibitem{bos1983}L. Bos, \emph{Bounding the Lebesgue function for Lagrange interpolation in a simplex}, J. Approx. Theory \textbf{38} (1983), 43--59. \url{https://doi.org/10.1016/0021-9045(83)90140-5}. The planar Fej\'er result is restated in Proposition 4.1 of\cite{bos2023}. The 1983 full text was not independently retrieved; this report also supplies its own elementary proof of the needed planar case.
\bibitem{muirhead}R. F. Muirhead, \emph{Some Methods applicable to Identities and Inequalities of Symmetric Algebraic Functions of n Letters}, Proc. Edinburgh Math. Soc. \textbf{21} (1902), 144--162. \url{https://doi.org/10.1017/S001309150003460X}.
\end{thebibliography}
\end{document}
